% !PRINTED 638
\weberpage{638}
%<<CONT
enthaltene Funktion, wenn die $x_1, x_2, \ldots, x_n$ ganze rationale
Funktionen sind, eine ganze Funktion in $\Omega$. Es ist aber nicht
gesagt, daß auch umgekehrt in der Form (2) mit \gesperrt{ganzen} Koeffizienten
 $x_i$ alle ganzen Funktionen in $\Omega$ enthalten sind.

Gibt es also ganze Funktionen in der Form (2), in der die
$x_i$ nicht alle ganze rationale Funktionen sind, so können wir eine
ganze Funktion finden:

\[
\frac{x_1 \omega_1 + x_2 \omega_2 + \cdots + x_n \omega_n}{z - c},
\]

in der die $x_i$ ganz und nicht alle durch $z - c$ teilbar sind.
Reduzieren wir die $x_i$ auf ihre konstanten Reste (nach $z - c$),
so ergibt sich eine ganze Funktion $\omega$ in der Form:

\weq{3}{\omega = \frac{c_1 \omega_1 + c_2 \omega_2 + \cdots + c_n \omega_n}{z - c},}

in der die Konstanten $c_1, c_2, \ldots, c_n$ jedenfalls nicht alle verschwinden.
Ist etwa $c_1$ von Null verschieden, so können wir $\omega_1$
durch $\omega, \omega_2, \ldots, \omega_n$ ausdrücken und erhalten eine neue ganzzahlige
 Basis von $\Omega$:

\weq{4}{\omega, \omega_2, \ldots, \omega_n.}

Für die Diskriminante ergibt sich aber nach § 173, (6) die
Beziehung:

\weq{5}{\varDelta(\omega, \omega_2, \ldots, \omega_n) = \frac{c_1^2}{(z-c)^2} \varDelta(\omega_1, \omega_2, \ldots, \omega_n).}

Beide Diskriminanten sind ganze rationale Funktionen von $z$,
aber die Diskriminante der Basis (4) ist von \gesperrt{niedrigerem Grade
als die Diskriminante der Basis} (1).

Wenn wir auf diese Weise fortfahren, den Grad der Diskriminante
 zu erniedrigen, müssen wir schließlich zu einer aus
ganzen Funktionen bestehenden Basis $\omega_1, \omega_2, \ldots, \omega_n$ kommen von
der Eigenschaft, daß in der Form (2) mit ganzen rationalen $x_i$
alle und nur die ganzen Funktionen in $\Omega$ ausgedrückt sind.

8. \gesperrt{Definition: Eine ganzzahlige Basis $\omega_1, \omega_2, \ldots, \omega_n$
des Körpers $\Omega$ heißt eine Minimalbasis, wenn in der
Form}

\weq{6}{\omega = x_1 \omega_1 + x_2 \omega_2 + \cdots + x_n \omega_n}

\gesperrt{mit ganzen rationalen $x_1, x_2, \ldots, x_n$ alle ganzen
Funktionen des Körpers $\Omega$ darstellbar sind.}

Eine solche Minimalbasis existiert also immer.
